\subsection{Cache 与访存层次设计}

意味核采用 Icache/Dcache + 统一 L2 + AXI 桥的层次结构，并在核内用 Store Queue 与 Store Buffer 分离推测写（STQ猜它会写，就是投机的意思）与提交写（SB确实要写）。

\subsubsection{一级指令缓存（ICache）}

ICache 容量 16\,KB，组织为 4 路$\times$128 组$\times$32\,B 行，采用 VIPT 查询，面向整行取指。
命中时向 IFU 返回一行指令数据；缺失时向 L2 请求填行。
支持 \texttt{cacop} 失效与 miss 途中的 \texttt{refill\_abort}：当同组维护请求到来时，丢弃迟到的安装，避免自修改场景下装入陈旧行。
替换在无效路优先策略下进行，存储使用 BRAM。

\subsubsection{一级数据缓存（DCache）}

DCache 与 ICache 相似（16\,KB，4 路$\times$128$\times$32\,B，VIPT），面向 LSU 的 load/store 访问。
具备非阻塞 miss 能力，MSHR 深度为 2；命中路径支持同行 store 合并，降低写后写/写后读气泡。
缺失时先查 L2；若需替换脏行，则按地址型维护把写回推进到 L2。
\texttt{cacop} 在完成 L1 写回/失效后，继续请求 L2 对同一物理行做写回并失效，以满足 Linux DMA 与缓存几何一致性。

\subsubsection{二级统一缓存（L2）}

L2 为 128\,KB 的 Icache 与 Dcache 共用的缓存，2 路$\times$2048 组$\times$32\,B，写回法。
Icache 与 Dcache 有 miss 时都会走 L2 ：Dcache 要 refill 时，支持数据边到边往 Dcache 送， load miss 能更快结束；I-miss 发现同一行上有 UC 写或 cache 维护写在进行时，不允许这次 I-miss 把那份旧行写回。
L2 miss 则经 \texttt{axi\_line\_bridge} 访问内存与外设。

\subsubsection{STQ 与 Store Buffer}

\begin{itemize}
  \item \textbf{STQ（深度 16）}：保存尚未提交的 store，维护顺序，提供 load 前递查询；指令提交后按 ROB 编号释放；
  \item \textbf{Store Buffer（深度 8）}：提交后写缓冲，可做行聚合与同字旁路合并，再写 Dcache/L2；冲刷时不清空，保证已提交的访存操作不丢失。
\end{itemize}

LSU 优先检查 SB/STQ 前递，再访问 Dcache ；命中可走旁路提前完成 ROB，提升访存密集的双提交率。
